% Reliability chapter (H5'). Numbers from results/final/unc_heldout.json and stats_held.json (H5 with image-bootstrap intervals); the
% exploratory recalibration (SAP_v3 H9) from paper/tables/v3_bone_intervals.tex (tools/hybrid_stats.py --tables_from).
\section{Design}
The deep ensemble of Chapter~\ref{ch:v2} gives, for every pixel, the mean $\bar f$ of the three bone estimates and their standard deviation
$\sigma$, the spread. Two questions follow (H5$'$). Does the spread rank the actual error, so that pixels where the members disagree are the
pixels where the bone estimate is wrong? And can the spread be turned into intervals that cover the true bone value at a stated rate? For the
second, a split-conformal scale $q$ is calibrated on the Qure-projection validation CTs, where the bone image is exact, so that
$[\bar f-q(\sigma+\varepsilon),\,\bar f+q(\sigma+\varepsilon)]$ has nominal coverage 0.90. Coverage is measured on the Qure-projection test
split (exact truth) and, as a measurement under domain shift, on the 51 held-out SZCH pairs against a pseudo-truth bone image: the radiograph
minus the dual-energy soft-tissue image after a least-squares gain and offset. The method is described in Section~\ref{sec:method_v2}.

\section{Results}
On the Qure-projection test split the intervals cover 0.900 $[0.897,0.903]$ of the pixels at the nominal 0.90, with a mean width of 0.018, and
the spread ranks the absolute error well (Spearman 0.76; Table~\ref{tab:reliability}). The calibration holds where it was made. On the held-out
SZCH pairs it does not: the same intervals cover 0.139 $[0.134,0.145]$ of the pixels, and the spread ranks the error only weakly (mean
per-image Spearman 0.11 $[0.07,0.15]$; Figure~\ref{fig:reliability}). Every one of the 51 images has coverage far below the nominal level, so
the failure is systematic rather than driven by a few images. H5$'$ holds on synthetic data and fails under the shift to real radiographs.
\begin{table}[htbp]
\centering\small
\caption{Ensemble spread and split-conformal intervals ($K=3$, nominal coverage 0.90), with image-bootstrap intervals.}
\label{tab:reliability}
\fittable{tables/reliability.tex}
\end{table}
\begin{figure}[htbp]
\centering
\safefig{figures/reliability.pdf}
\caption{Left: per-image coverage of the conformal intervals on the Qure-projection test split and on the held-out SZCH-X-Rays pairs (bars:
means; dashed: nominal 0.90). Right: per-image Spearman correlation of ensemble spread and absolute error on SZCH-X-Rays (line: mean).}
\label{fig:reliability}
\end{figure}

\section{Recalibration on real pairs (exploratory)}
After this failure was known, the third analysis plan (H9, Chapter~\ref{ch:classifier}) added a post hoc check on the same, already-seen
pairs: the intervals were recalibrated on the 24 SZCH development pairs and tested on the 51 held-out pairs. One variant uses a single global
calibration; a Mondrian variant calibrates separately per stratum of predicted bone overlap and spread tercile. Pooled coverage was 0.886
with the global calibration and 0.895 with the Mondrian calibration, whose intervals were less than half as wide (mean width 0.215 against
0.447; Table~\ref{tab:bone_recal}); the median image reached 0.886 and 0.894. Figure~\ref{fig:bone_recal} also shows a per-image variant that
takes the calibration quantile over per-image quantiles. A few real pairs therefore restore nominal coverage on average. The check is
exploratory: it was designed after the failure was seen, it uses pairs from one hospital, and the coverage of single images still varies.
\begin{table}[htbp]
\centering\small
\caption{Exploratory recalibration on real pairs (post hoc): bone-image intervals calibrated on the 24 SZCH-X-Rays development pairs and tested
on the 51 held-out pairs at nominal 0.90. Mondrian: calibration per stratum of predicted bone overlap and ensemble spread.}
\label{tab:bone_recal}
\fittable{tables/v3_bone_intervals.tex}
\end{table}
\begin{figure}[htbp]
\centering
\safefig{figures/fig_bone_uncertainty.pdf}
\caption{Exploratory recalibration on real pairs (post hoc). (a) Per-image coverage on the 51 held-out SZCH-X-Rays pairs for global and
Mondrian calibration, with quantiles pooled over pixels or taken per image; dashed: nominal 0.90. (b) Reliability of the Mondrian intervals
across nominal levels; cross: the synthetic calibration. (c) Sparsification of the absolute bone error by the ensemble spread, with the
error-sorted oracle.}
\label{fig:bone_recal}
\end{figure}

\section{Interpretation and limits}
Three things explain why the intervals fail on real images, and they cannot be separated with the data at hand. First, the spread rises only
modestly on real radiographs (mean interval width 0.026 against 0.018) while the error rises far more: the members agree with each other more
than they agree with the dual-energy reference, because they share their training data and their biases. Second, the pseudo-truth bone image is not the true bone image;
it carries the differences in tone and texture between a radiograph and a dual-energy soft-tissue image, so part of the measured error is
reference error. Third, split-conformal calibration guarantees coverage only for data exchangeable with the calibration data, and only marginally
over pixels, not per image; rendered projections and real radiographs are not exchangeable. The spread is still useful as a flag: on the
paediatric radiographs of Chapter~\ref{ch:external} it is 2.7 times its value on adult pairs. Calibrated intervals on real radiographs need
calibration data with real truth; the exploratory recalibration above shows that 24 dual-energy pairs from one site were enough there, not
whether they would transfer to other sites.
